\section{Introduction}

EEG brain--computer interfaces must generalize across subjects, yet inter-subject variability is large and
a decoder can latch onto subject identity rather than the neural task. A growing line of work treats this
as a domain-generalization (DG) problem and as a \emph{shortcut/leakage} problem: features that separate
subjects, conditioned on the label, are a direct measure of how much the representation could be exploiting
domain identity. Graph neural networks are attractive for EEG because electrodes form a natural graph, and
they expose interpretable objects --- a graph-level readout, per-electrode node features, and (for some
designs) a learned adjacency. This makes them a good substrate for \emph{auditing} leakage at the graph and
node level.

Most leakage-control work either presumes a task-capable model or measures leakage without first confirming
that the backbone can learn the task at all; on EEG graph networks, we find neither holds for free. This
paper therefore makes a deliberately \textbf{bounded} contribution. We \emph{audit} label-conditional
domain leakage in a graph network's \texttt{graph\_z}/\texttt{node\_z} objects under a strict
no-target-label firewall; we then show --- through a methodology chain whose negative results are part of
the evidence --- that leakage control is only meaningful on a \emph{task-capable} graph backbone, which is
nontrivial to obtain; and we finally confirm that a single fixed graph/node conditional-information
regularizer reduces the audited leakage on two MI datasets while meeting a pre-registered source-task
retention gate. We do not claim state-of-the-art accuracy, leakage elimination, an unbiased CMI estimate,
an edge-level method, or generality beyond this backbone and these two MI datasets. The contributions are:

\begin{itemize}
  \item[\textbf{(C1)}] A \textbf{source-only audit} of label-conditional domain leakage in a graph
  network's \texttt{graph\_z}/\texttt{node\_z} objects: a posterior-KL plug-in proxy with a retrained
  within-label permutation null and a target-label firewall (we explicitly disclaim unbiased CMI).
  \item[\textbf{(C2)}] A \textbf{methodology chain, including its negative results}, showing that a
  \emph{task-capable} graph backbone is a prerequisite --- the original graph backbone is near-chance,
  dynamic-edge designs overfit, and a static-adjacency DGCNN is the one that learns the task.
  \item[\textbf{(C3)}] A \textbf{fixed-candidate, cross-dataset confirmation} that a graph/node
  conditional-information regularizer ($\lambda_g = \lambda_n = 0.010$, no edge term) partially and
  reproducibly reduces the audited leakage while meeting the pre-registered source-task retention gate on
  two MI datasets (the gate holds dataset-wide, with one BNCI2015\_001 fold missing the per-fold threshold).
  \item[\textbf{(C4)}] An explicit \textbf{negative-results $+$ limitations} account that bounds the claim
  (proxy not unbiased CMI; partial not full; graph/node only; one backbone; two MI datasets).
\end{itemize}
